% انتخاب کلاس سند
\documentclass[12pt,a4paper]{article}

% ---------- پکیج‌ها (همه قبل از xepersian) ----------
\usepackage{fontspec}
\usepackage{amsmath, amssymb, amstext}
\usepackage{graphicx}
\usepackage{geometry}
\usepackage{xcolor}
\usepackage{fancyhdr}

% ---------- تنظیم حاشیه ----------
\geometry{margin=1in}

% ---------- تنظیم fancyhdr ----------
\pagestyle{fancy}
\setlength{\headheight}{15pt} % حل خطای headheight

\lhead{کوییز درس مکانیک تحلیلی}
\chead{دانشکده فیزیک دانشگاه صنعتی خواجه نصیرالدین طوسی}
\rhead{\today}

% ---------- فارسی (باید آخر بیاید) ----------
\usepackage{xepersian}
\settextfont{B Nazanin}

% ---------- شروع سند ----------
\begin{document}
	
	\begin{tabular}{|p{7cm}|p{7cm}|}
		\hline
		\textbf{نام و نام خانوادگی:} & \textbf{شماره دانشجویی:} \\
		\hline
	\end{tabular}
	
	\section*{سوالات:}
	
	\begin{enumerate}
		\item 
		معادلات حرکت پرتابه را در میدان گرانشی یکنواخت و بدون مقاومت هوا بدست آورید
		\item 
			ذره ای در امتداد یک جوی چرخزاد صاف آزادانه حرکت میکند.سطح این چرخزاد در قالب این معادلات پارامتری بیان میشود
			\\
	\[
	x=\frac{a}{4}\bigl(2\theta+\sin\theta\bigr)
	\]
	\[
	y=\frac{a}{4}\bigl(1-\cos(2\theta)\bigr)
	\]که در آن $\theta \in [0,\pi]$ و $a$ ثابت است.
	لانگرانژی و معادلات حرکت ذره را بیابید.
	
	\end{enumerate}
	
	\vspace{10pt}
	\textbf{موفق باشید.}
	
\end{document}
